\section{Conclusion}
\label{sec:conclusion}

Industrial automation spent three decades building machine-readable meaning for
a purpose that has now changed. The standards were written so that two machines
could agree; they turn out to be the scaffolding that lets a language model be
\emph{right on purpose} rather than by coincidence.

This survey mapped that landscape through an agent's eyes and then tested the map.
The map says that ``semantics'' decomposes into ten separable capabilities, and
that industrial standardisation has solved the typing problem far more completely
than the unit, behaviour, constraint and provenance problems. The test says that
this imbalance is expensive. A fully typed information model, the state of
practice and the endpoint of most industrial AI programmes, supports
\CeilingCTwo{} of our task suite, catches \RecallGTwo{} of faulty agent actions,
raises false alarms on \FPRGTwo{} of legal ones, forbids three quarters of the
legal action space, and does not improve when given a context window
$32\times$ larger. Adding a unit calculus, a plant graph, a behavioural model,
an explicit interlock model and observation provenance takes the same suite to
\CeilingCFour{}, guardrail recall to \RecallGFour{} with no false alarms, and
action-space precision to \ActPrecGFour{} at full recall.

Three recommendations follow, in order of return on effort. \emph{Adopt a unit
ontology}: it is small, mechanical, and it eliminates both a class of silent
errors and a class of false alarms. \emph{Publish the behaviour and the
constraints}, namely the PackML states, the interlock guards and the operating envelopes,
as models rather than as ladder logic and PDFs. They already exist, and they
are the only thing that catches the failures that matter. \emph{Put the semantic
layer behind a query interface and a validator}, not into the prompt: serialised
industrial models cost up to $\MaxSerialRatio\times$ the tokens of a tag list and
still cannot answer a plant-wide question that a $222$-token query answers
exactly.

Two objections deserve the answers we measured rather than the answers we
assumed. The first is that tag names already carry the semantics, so none of
this is needed. They partly do, a resolver that decodes mnemonics reaches
\ConvMnemonicDecoded{} on the convention it was built for, but the capability
swings by \ConvSpreadCOnePlusconv{} points across plants and, where a
convention has stopped being truthful, produces confident errors at
\ConvLyingWrong{}. Semantics carried by a convention is semantics on which an
agent cannot be corrected. The second is that richer models simply cost more
context than they are worth. At a fixed budget they do, until the retrieval
layer is told to fetch the semantic layer a question actually needs, at which
point the ceiling rises from \RouteCFourNeighCeiling{} to
\RouteCFourRoutedCeiling{} at identical cost. Richness is not expensive;
projecting it badly is.

The broader point is methodological. Debate about agent reliability tends to
focus on the model, its scale, its reasoning, its prompt. In an industrial
setting a large part of the reliability is decided before the model is invoked,
by whether the representation makes the right answer \emph{expressible} and the
wrong answer \emph{refutable}. That is a property of the data layer, it is
measurable. As the ceiling curves show, no amount of model capability
compensates for getting it wrong.

Every number in this paper is reproducible offline with one command. We hope the
benchmark, the fault taxonomy and the agent-usability framework are more useful
than the specific figures, which belong to one synthetic plant. The figures will
change; the ordering, we expect, will not.
